\documentclass[11pt]{article}

\usepackage{amsmath,amssymb,amsthm}
\usepackage{geometry}
\usepackage{url}

\geometry{margin=1in}

\title{\textbf{Prefactor Coefficient Laws\\
for the Gauss--Legendre AGM Algorithm}}

\author{Wayne Baker\\
\small Independent Researcher\\
\small Goulds, Newfoundland and Labrador, Canada}

\date{September 21, 2026}

\theoremstyle{plain}
\newtheorem{theorem}{Theorem}
\newtheorem{lemma}{Lemma}

\theoremstyle{definition}
\newtheorem{definition}{Definition}
\newtheorem{remark}{Remark}

\begin{document}

\maketitle

\begin{abstract}
We study two fitted coefficients associated with the Gauss--Legendre arithmetic--geometric mean (AGM) algorithm for $\pi$. For consecutive Gauss--Legendre approximants $\pi_n$, define
\[
\alpha_n
=
-\frac{\pi_{n+1}-\pi}{\pi_{n+1}-\pi_n}.
\]
This coefficient is the exact one--step weight which corrects $\pi_{n+1}$ to
$\pi$ using the increment $\pi_{n+1}-\pi_n$.

The first result is a coefficient law for the Gauss--Legendre error operator:
\[
\frac{\alpha_{n+1}}{\alpha_n^2}
\longrightarrow
\frac12.
\]
The proof uses Brent's theta--function expansion for the lower
Gauss--Legendre approximant. In the notation of the present paper, that expansion yields the asymptotic
\[
\pi-\pi_n
=
\left(2^{n+4}\pi^2-8\pi\right)e^{-\pi 2^{n+1}}
+
O\!\left(2^n e^{-\pi 2^{n+2}}\right).
\]
Consequently, if
\[
\varepsilon_n=\pi_n-\pi=A_n e^{-\pi 2^{n+1}},
\]
then
\[
A_n\sim -2^{n+4}\pi^2,
\]
and hence $A_{n+1}/A_n\to 2$. The limiting coefficient $1/2$ thus arises from
exact cancellation of the doubly--exponential factor in the quotient
$\alpha_{n+1}/\alpha_n^2$, leaving a prefactor ratio that tends to $1/2$.

For the companion numerator--channel coefficients $\beta_n$, a direct
theta--function analysis gives the sharper prefactor law
\[
\beta_n
=
-\left(2^{n+2}\pi-1\right)e^{-\pi 2^{n+1}}
\left(1+O\!\left(e^{-\pi 2^{n+1}}\right)\right),
\]
and consequently
\[
-\pi 2^{n+1}\frac{\beta_{n+1}}{\beta_n^2}\longrightarrow1.
\]
Thus the previously observed decay $\beta_{n+1}/\beta_n^2\to0$ is analytic,
with an explicit leading scale. Finally, we explain why the static Baker
N--series cancellation operator does not transfer directly to the
Gauss--Legendre setting, while the natural dynamic cancellation principle is
closely related to Aitken's $\Delta^2$ process.
\end{abstract}

\section{Preliminaries}

Let $\{\pi_n\}$ denote the Gauss--Legendre $\pi$--sequence generated by the
AGM iteration
\[
\begin{aligned}
a_{n+1}&=\tfrac12(a_n+b_n), & b_{n+1}&=\sqrt{a_n b_n},\\
t_{n+1}&=t_n-p_n(a_n-a_{n+1})^2, & p_{n+1}&=2p_n,
\end{aligned}
\]
with initial values $a_0=1$, $b_0=1/\sqrt2$, $t_0=1/4$, $p_0=1$. The Gauss--Legendre
approximant is
\[
\pi_n=\frac{(a_n+b_n)^2}{4t_n}=\frac{a_{n+1}^2}{t_n}.
\]
Write $\varepsilon_n=\pi_n-\pi$ for the signed error. For the lower
Gauss--Legendre approximant considered here one has $\pi_n<\pi$ and hence
$\varepsilon_n<0$.

The Gauss--Legendre algorithm is classically known to converge extremely fast
(the number of correct digits roughly doubles per iteration). A more precise
description of the error is doubly exponential. Brent's theta--function analysis gives a sharp expansion for the lower Gauss--Legendre approximant; see Brent~\cite[Section~4, Eq.~(20),
pp.~9--10]{brent2020}:
\[
\pi-\pi_n
=
\left(2^{n+4}\pi^2-8\pi\right)e^{-\pi 2^{n+1}}
+
O\!\left(2^n e^{-\pi 2^{n+2}}\right).
\]

\begin{remark}
The phrase ``quadratic convergence'' should not be confused with a fixed
constant-$C$ model for every associated error sequence.  For Brent's upper
approximants $a_n^2/t_n$, the corresponding errors $e_n=a_n^2/t_n-\pi$
satisfy
\[
\frac{e_{n+1}}{e_n^2}\longrightarrow\frac1{8\pi};
\]
see Brent~\cite[Section~4, following Eq.~(19)]{brent2020}.  By contrast, for
the lower approximants used here,
$\varepsilon_n=\pi_n-\pi$ satisfies
$\varepsilon_{n+1}/\varepsilon_n^2\to0$.  Thus the lower channel has a
rapidly varying prefactor even though the overall iteration remains
quadratically convergent in the usual digits-doubling sense.
\end{remark}

\section{Fitted Operator Coefficients}

\begin{definition}
Define
\[
\alpha_n
=
-\frac{\pi_{n+1}-\pi}{\pi_{n+1}-\pi_n}
=
-\frac{\varepsilon_{n+1}}{\varepsilon_{n+1}-\varepsilon_n},
\]
so that the exact one-step identity
\[
\pi=\pi_{n+1}+\alpha_n(\pi_{n+1}-\pi_n)
\]
holds for each $n$.
\end{definition}

Thus $\alpha_n$ is the exact fitted coefficient which recovers $\pi$ from the
two successive approximants.

For later comparison, define the numerator sequence $Q_n=(a_n+b_n)^2$ and the
numerator-channel coefficient
\[
\beta_n=-\frac{Q_{n+1}-4\pi t_{n+1}}{Q_{n+1}-Q_n},
\]
which plays an analogous role on the numerator channel.

\section{Reduction to Error Ratios}

\begin{lemma}\label{lem:alpha}
Assume $r_n:=\varepsilon_{n+1}/\varepsilon_n\to 0$. Then
\[
\alpha_n\sim \frac{\varepsilon_{n+1}}{\varepsilon_n}=r_n,
\]
and
\[
\frac{\alpha_{n+1}}{\alpha_n^2}\sim \frac{r_{n+1}}{r_n^2}
=\frac{\varepsilon_{n+2}\varepsilon_n^2}{\varepsilon_{n+1}^3}.
\]
\end{lemma}

\begin{proof}
From $\alpha_n=\dfrac{r_n}{1-r_n}$ and $r_n\to0$ we obtain
$\alpha_n\sim r_n$. More precisely,
\[
\frac{\alpha_{n+1}}{\alpha_n^2}
=
\frac{r_{n+1}}{r_n^2}
\frac{(1-r_n)^2}{1-r_{n+1}}.
\]
Since $r_n,r_{n+1}\to0$, the correction factor tends to $1$, and therefore
\[
\frac{\alpha_{n+1}}{\alpha_n^2}
\sim
\frac{r_{n+1}}{r_n^2}
=
\frac{\varepsilon_{n+2}\varepsilon_n^2}{\varepsilon_{n+1}^3}.
\]
\end{proof}

\section{Prefactor Asymptotic}

\begin{lemma}\label{lem:prefactor}
For the lower Gauss--Legendre approximants one has
\[
\varepsilon_n=\pi_n-\pi
=
-\left(2^{n+4}\pi^2-8\pi\right)e^{-\pi 2^{n+1}}+O\!\left(2^n e^{-\pi 2^{n+2}}\right).
\]
Consequently, writing $\varepsilon_n=A_n e^{-\pi 2^{n+1}}$, we have
\[
A_n=-\left(2^{n+4}\pi^2-8\pi\right)+O\!\left(2^n e^{-\pi 2^{n+1}}\right),
\]
so $A_n\sim -2^{n+4}\pi^2$ and hence $A_{n+1}/A_n\to2$.
\end{lemma}

\begin{proof}
Since
\[
a_{n+1}
=
\frac{a_n+b_n}{2},
\]
the approximant used in this paper satisfies
\[
\pi_n
=
\frac{(a_n+b_n)^2}{4t_n}
=
\frac{a_{n+1}^2}{t_n}.
\]
Brent writes the same Gauss--Legendre auxiliary sequence as $s_n$ rather than
$t_n$, so the present approximant is Brent's lower Gauss--Legendre approximant
\[
\frac{a_{n+1}^2}{s_n}.
\]
Brent's theta--function expansion for this lower approximant gives
\[
\pi-\frac{a_{n+1}^2}{s_n}
=
\left(2^{n+4}\pi^2-8\pi\right)e^{-\pi 2^{n+1}}
+
O\!\left(2^n e^{-\pi 2^{n+2}}\right);
\]
see Brent~\cite[Section~4, Eq.~(20), pp.~9--10]{brent2020}.
Therefore
\[
\pi-\pi_n
=
\left(2^{n+4}\pi^2-8\pi\right)e^{-\pi 2^{n+1}}
+
O\!\left(2^n e^{-\pi 2^{n+2}}\right).
\]
Since
\[
\varepsilon_n
=
\pi_n-\pi,
\]
we obtain
\[
\varepsilon_n
=
-\left(2^{n+4}\pi^2-8\pi\right)e^{-\pi 2^{n+1}}
+
O\!\left(2^n e^{-\pi 2^{n+2}}\right).
\]
Writing
\[
\varepsilon_n
=
A_n e^{-\pi 2^{n+1}},
\]
gives
\[
A_n
=
-\left(2^{n+4}\pi^2-8\pi\right)
+
O\!\left(2^n e^{-\pi 2^{n+1}}\right).
\]
Hence
\[
A_n
\sim
-2^{n+4}\pi^2,
\]
and therefore
\[
\frac{A_{n+1}}{A_n}
\to
2.
\]
\end{proof}

\section{The $\alpha_n$ Coefficient Law}

\begin{theorem}\label{thm:main}
With $\alpha_n$ as above, one has
\[
\lim_{n\to\infty}\frac{\alpha_{n+1}}{\alpha_n^2}=\frac12.
\]
\end{theorem}

\begin{proof}
By Lemma~\ref{lem:prefactor} we may write
$\varepsilon_n=A_n e^{-\pi 2^{n+1}}$ with $A_{n+1}/A_n\to2$. Hence
$\varepsilon_{n+1}/\varepsilon_n\to0$, so Lemma~\ref{lem:alpha} applies and
\[
\frac{\alpha_{n+1}}{\alpha_n^2}\sim\frac{\varepsilon_{n+2}\varepsilon_n^2}{\varepsilon_{n+1}^3}.
\]
Inserting the asymptotic form cancels the exponential contributions
exactly (the exponent combination equals zero), leaving
\[
\frac{\varepsilon_{n+2}\varepsilon_n^2}{\varepsilon_{n+1}^3}\sim
\frac{A_{n+2}A_n^2}{A_{n+1}^3}
=
\frac{A_{n+2}}{A_{n+1}}\left(\frac{A_n}{A_{n+1}}\right)^2.
\]
Using $A_{n+1}/A_n\to2$ we get
\[
\frac{A_{n+2}}{A_{n+1}}\to 2,\qquad \frac{A_n}{A_{n+1}}\to\tfrac12,
\]
and therefore the prefactor quotient tends to
$2\cdot(\tfrac12)^2=\tfrac12$. The theorem follows.
\end{proof}

\begin{remark}
The limit $1/2$ is a prefactor fingerprint of the AGM structure: it comes
from exact cancellation of the doubly--exponential factor together with the
linear--in--$2^n$ growth of $A_n$. It does not follow from a fixed quadratic
recurrence $\varepsilon_{n+1}\sim C\varepsilon_n^2$.
\end{remark}

\section{Numerical Verification of the $\alpha_n$ Law}

All reported numerical values were computed with \texttt{mpmath} at
500 decimal digits, reconstructing the Gauss--Legendre sequence from its
defining recursion and using \texttt{mp.pi} for the reference value of
$\pi$. The fitted coefficients and diagnostic ratios support the theoretical limit
over the precision and ranges shown in Tables~\ref{tab:alpha} and~\ref{tab:ratio}.
As an independent check, the direct error quotient
$\varepsilon_{n+2}\varepsilon_n^2/\varepsilon_{n+1}^3$ approaches
$\alpha_{n+1}/\alpha_n^2$ rapidly.  The two quantities are not identical at
finite $n$; their exact relation is
\[
\frac{\alpha_{n+1}}{\alpha_n^2}
=
\frac{\varepsilon_{n+2}\varepsilon_n^2}{\varepsilon_{n+1}^3}
\frac{(1-r_n)^2}{1-r_{n+1}},
\qquad
r_n=\frac{\varepsilon_{n+1}}{\varepsilon_n}.
\]
The correction factor tends to $1$, explaining the increasing numerical
agreement in Table~\ref{tab:ratio}.

\begin{table}[h]
\centering
\begin{tabular}{c l}
\hline
$n$ & $\alpha_n$ \\
\hline
0 & $4.476840443051657\times 10^{-3}$ \\
1 & $7.278747106485182\times 10^{-6}$ \\
2 & $2.482705775153228\times 10^{-11}$ \\
3 & $2.988090953871379\times 10^{-22}$ \\
4 & $4.397065311352923\times 10^{-44}$ \\
5 & $9.594611535813435\times 10^{-88}$ \\
6 & $4.585616315217315\times 10^{-175}$ \\
\hline
\end{tabular}
\caption{Fitted coefficients $\alpha_n$.}
\label{tab:alpha}
\end{table}

\begin{table}[h]
\centering
\begin{tabular}{c l l}
\hline
$n$ & $\alpha_{n+1}/\alpha_n^2$ & $A_{n+1}/A_n$ \\
\hline
0 & $0.3631728700311416$ & $2.3866261558728965$ \\
1 & $0.4686104763344650$ & $2.0871750987545307$ \\
2 & $0.4847784478028806$ & $2.0414374856886031$ \\
3 & $0.4924649101466708$ & $2.0202981874959604$ \\
4 & $0.4962511920458477$ & $2.0100471245391349$ \\
5 & $0.4981302572056423$ & $2.0049984522335209$ \\
\hline
\end{tabular}
\caption{Coefficient and prefactor ratios: the left column approaches $1/2$
from below; the right column approaches $2$, consistent with the prefactor
asymptotic.}
\label{tab:ratio}
\end{table}

\section{The Numerator-Channel Coefficients}

For this section let Brent's elliptic nome be
\[
q=e^{-\pi},
\]
and define
\[
x_n=q^{2^{n+1}}=e^{-\pi 2^{n+1}},
\qquad
C_n=2^{n+2}\pi-1.
\]
Also set
\[
M=\operatorname{AGM}\!\left(1,\frac1{\sqrt2}\right)
=\theta_3^{-2}(q).
\]
The capital notation \(Q_n=(a_n+b_n)^2\) distinguishes the numerator sequence
from the nome \(q\).

\begin{lemma}\label{lem:beta-theta}
As \(n\to\infty\),
\[
Q_{n+1}-Q_n
=
-32M^2x_n
\left(
1+2x_n+4x_n^2+O(x_n^4)
\right),
\]
and
\[
Q_{n+1}-4\pi t_{n+1}
=
-32M^2C_nx_n^2
\left[
1+\left(2-\frac1{C_n}\right)x_n^2+O(x_n^4)
\right].
\]
The implied constants may be chosen independently of \(n\).
\end{lemma}

\begin{proof}
Brent's theta parametrization gives
\[
a_n
=
\frac{\theta_3^2(q^{2^n})}{\theta_3^2(q)}
=
M\theta_3^2(q^{2^n});
\]
see Brent~\cite[Section~4, Eq.~(13)]{brent2020}. Since
\(a_{n+1}=(a_n+b_n)/2\),
\[
Q_n=4a_{n+1}^2=4M^2\theta_3^4(x_n).
\]
Because \(x_{n+1}=x_n^2\),
\[
Q_{n+1}-Q_n
=
4M^2\left[\theta_3^4(x_n^2)-\theta_3^4(x_n)\right].
\]
From the defining theta series,
\[
\theta_3(x)=1+2x+2x^4+O(x^9),
\]
hence
\[
\theta_3^4(x)
=
1+8x+24x^2+32x^3+24x^4+48x^5+O(x^6).
\]
It follows that
\[
\theta_3^4(x^2)-\theta_3^4(x)
=
-8x-16x^2-32x^3+O(x^5),
\]
which proves the first expansion.

For the second difference, Brent's exact limit and tail formulas are
\[
t_\infty=\frac{M^2}{\pi},
\qquad
t_{n+1}-t_\infty
=
M^2\sum_{m=n+1}^{\infty}
2^m\theta_2^4(q^{2^{m+1}});
\]
see Brent~\cite[Section~4, Eqs.~(17)--(18)]{brent2020}. Therefore
\[
Q_{n+1}-4\pi t_{n+1}
=
4M^2
\left[
\theta_3^4(x_n^2)-1
-
\pi\sum_{m=n+1}^{\infty}
2^m\theta_2^4(q^{2^{m+1}})
\right].
\]
The defining series gives
\[
\theta_2(y)
=
2y^{1/4}\left(1+y^2+O(y^6)\right),
\qquad
\theta_2^4(y)
=
16y+64y^3+O(y^5).
\]
Reindexing the tail with \(m=n+1+j\) yields
\[
\sum_{m=n+1}^{\infty}
2^m\theta_2^4(q^{2^{m+1}})
=
2^{n+5}x_n^2
+
2^{n+6}x_n^4
+
O(2^n x_n^6),
\]
where the terms \(j\ge2\) are smaller still. Also
\[
\theta_3^4(x_n^2)-1
=
8x_n^2+24x_n^4+O(x_n^6).
\]
Thus
\begin{align*}
\frac{Q_{n+1}-4\pi t_{n+1}}{4M^2}
&=
(8-2^{n+5}\pi)x_n^2
+
(24-2^{n+6}\pi)x_n^4
+
O(2^n x_n^6)\\
&=
-8C_nx_n^2
-
8(2C_n-1)x_n^4
+
O(2^n x_n^6).
\end{align*}
Since \(C_n\asymp2^n\), division by the leading term gives the stated second
expansion.
\end{proof}

\begin{theorem}[Numerator-channel prefactor law]
\label{thm:beta-prefactor}
The numerator--channel coefficient satisfies
\[
\boxed{
\beta_n
=
-C_nx_n
\left[
1-2x_n+
\left(2-\frac1{C_n}\right)x_n^2
+O(x_n^3)
\right].
}
\]
In particular,
\[
\beta_n
\sim
-\left(2^{n+2}\pi-1\right)e^{-\pi2^{n+1}}.
\]
Moreover,
\[
\boxed{
\frac{\beta_{n+1}}{\beta_n^2}
=
-\frac{C_{n+1}}{C_n^2}
\left[
1+4x_n+
\left(6+\frac{2}{C_n}\right)x_n^2
+O(x_n^3)
\right].
}
\]
Consequently,
\[
\frac{\beta_{n+1}}{\beta_n^2}
\sim
-\frac{1}{\pi2^{n+1}},
\qquad
\boxed{
-\pi2^{n+1}\frac{\beta_{n+1}}{\beta_n^2}\longrightarrow1.
}
\]
\end{theorem}

\begin{proof}
By Lemma~\ref{lem:beta-theta},
\[
\beta_n
=
-C_nx_n
\frac{
1+\left(2-C_n^{-1}\right)x_n^2+O(x_n^4)
}{
1+2x_n+4x_n^2+O(x_n^4)
}.
\]
Since
\[
\frac1{1+2x+4x^2+O(x^4)}
=
1-2x+8x^3+O(x^4),
\]
multiplication gives the first expansion.

Now \(x_{n+1}=x_n^2\). Applying the first expansion at \(n\) and \(n+1\)
gives
\[
\beta_{n+1}
=
-C_{n+1}x_n^2
\left[1-2x_n^2+O(x_n^4)\right]
\]
and
\[
\beta_n^2
=
C_n^2x_n^2
\left[
1-4x_n+
\left(8-\frac{2}{C_n}\right)x_n^2
+O(x_n^3)
\right].
\]
Taking the quotient and expanding the reciprocal yields the stated quotient
formula. Finally,
\[
\frac{C_{n+1}}{C_n^2}
=
\frac{2^{n+3}\pi-1}{(2^{n+2}\pi-1)^2}
\sim
\frac{1}{\pi2^{n+1}},
\]
which proves the scaled limit.
\end{proof}

\begin{remark}
The theorem explains the formerly numerical observation that
\(\beta_{n+1}/\beta_n^2\) tends to zero with an approximately halving
finite-\(n\) scale. To avoid overloading the prefactor notation \(A_n\)
used earlier, set
\[
D_n=2^{n+2}\pi.
\]
Then
\[
\pi2^{n+1}\frac{C_{n+1}}{C_n^2}
=
\frac{D_n(2D_n-1)}{2(D_n-1)^2}
=
1+\frac{3}{2D_n}+\frac{2}{D_n^2}+O(D_n^{-3}).
\]
Thus
\[
-\pi2^{n+1}\frac{\beta_{n+1}}{\beta_n^2}
=
1+\frac{3}{2^{n+3}\pi}
+O(4^{-n})
+O(x_n).
\]
\end{remark}

\section{Numerical Verification of the $\beta_n$ Law}

The companion verification script evaluates the numerator--channel
coefficient directly from the Gauss--Legendre recursion at high precision.
Table~\ref{tab:beta} records the coefficient itself together with the scaled
quantity
\[
S_n
=
-\pi2^{n+1}\frac{\beta_{n+1}}{\beta_n^2},
\]
which Theorem~\ref{thm:beta-prefactor} proves tends to $1$.

\begin{table}[ht]
\centering
\begin{tabular}{c l l}
\hline
$n$ & $\beta_n$ & $S_n$ \\
\hline
0 & $-2.151900766752290\times10^{-2}$ & $1.141914530829174$ \\
1 & $-8.415854367694321\times10^{-5}$ & $1.063029581987583$ \\
2 & $-5.991449587079527\times10^{-10}$ & $1.030653289501889$ \\
3 & $-1.472097424691757\times10^{-20}$ & $1.015121159164455$ \\
4 & $-4.376441532326746\times10^{-42}$ & $1.007510170612647$ \\
5 & $-1.919516499821586\times10^{-85}$ & $1.003742600849157$ \\
\hline
\end{tabular}
\caption{Direct numerical check of the numerator--channel law.  The scaled
quantity $S_n$ approaches $1$ as predicted by
Theorem~\ref{thm:beta-prefactor}.}
\label{tab:beta}
\end{table}

The same run also checks the sharper normalized expansion
\[
\frac{\beta_n}{-C_nx_n}
=
1-2x_n+
\left(2-\frac1{C_n}\right)x_n^2
+O(x_n^3),
\]
with the displayed residuals decreasing doubly exponentially over the
resolved range.  These computations are verification only; the theorem is
proved analytically above.

\section{Relation to Static and Dynamic Cancellation Operators}

The static N--series cancellation operator used in
Baker~\cite{baker2026nseries} targets even--power expansions in a scale
parameter and uses fixed rational weights to cancel terms. The
Gauss--Legendre error is doubly exponential with a rapidly varying prefactor,
so no fixed rational weight yields exact cancellation across iterates. The
dynamic fitted coefficient $\alpha_n$ plays the role of a local, data--driven
cancellation weight, analogous in spirit to Aitken's $\Delta^2$ process which
eliminates a single geometric error mode estimated from consecutive iterates.

\section{Summary}

\begin{enumerate}
\item Brent's theta--function expansion yields the prefactor asymptotic
$A_n\sim -2^{n+4}\pi^2$ for the lower Gauss--Legendre approximants.
\item Consequently $A_{n+1}/A_n\to2$, and the fitted coefficients satisfy
$\alpha_{n+1}/\alpha_n^2\to 1/2$.
\item The limit $1/2$ arises from exact cancellation of the doubly--exponential
factor and from the prefactor quotient; it is not a consequence of a fixed
quadratic error map.
\item The numerator--channel coefficients satisfy
\[
\beta_n\sim-(2^{n+2}\pi-1)e^{-\pi2^{n+1}},
\qquad
-\pi2^{n+1}\frac{\beta_{n+1}}{\beta_n^2}\to1.
\]
\end{enumerate}

\paragraph{Reproducibility.}
Numerical values were generated by reconstructing the Gauss--Legendre sequence
at 500 decimal digits and computing the diagnostics
$\varepsilon_n,\alpha_n,\alpha_{n+1}/\alpha_n^2,
\varepsilon_{n+2}\varepsilon_n^2/\varepsilon_{n+1}^3$, and
$A_{n+1}/A_n$ with \texttt{mpmath}. A companion high--precision audit checks
the $\beta_n$ prefactor expansion and scaled quotient law. The reference $\pi$
was \texttt{mp.pi}.

\begin{thebibliography}{9}
\bibitem{aitken1926}
A.~C.~Aitken,
\emph{On Bernoulli's numerical solution of algebraic equations},
Proc. Roy. Soc. Edinburgh \textbf{46} (1926), 289--305.

\bibitem{borwein1987}
J.~M.~Borwein and P.~B.~Borwein,
\emph{Pi and the AGM: A Study in Analytic Number Theory and Computational Complexity},
John Wiley \& Sons, New York, 1987.

\bibitem{baker2026nseries}
W.~Baker,
\emph{A Constructive N--Series Acceleration Law for Polygonal Approximation of $\pi$},
preprint and reproducibility archive, 2026.

\bibitem{brent1976}
R.~P.~Brent,
\emph{Fast multiple--precision evaluation of elementary functions},
J. ACM \textbf{23} (1976), no.~2, 242--251.

\bibitem{brent2020}
R.~P.~Brent,
\emph{The Borwein brothers, Pi and the AGM},
in \emph{From Analysis to Visualisation}, Springer Proc. Math. Stat.
\textbf{313}, Springer, 2020, 323--348. Also available as arXiv:1802.07558.

\bibitem{salamin1976}
E.~Salamin,
\emph{Computation of $\pi$ using arithmetic--geometric mean},
Math. Comp. \textbf{30} (1976), no.~135, 565--570.
\end{thebibliography}

\end{document}